% Use only LaTeX2e, calling the article.cls class and 12-point type.
\RequirePackage{lineno}
\documentclass[11pt]{article}

\usepackage{lipsum}
\usepackage{amsmath,amssymb,cases,multirow, graphicx}
\usepackage{dcolumn, rotating}

%~~~~~~~~~~~
% Reasonable page setup
%~~~~~~~~~~~
\topmargin 0.0cm
\oddsidemargin 0.2cm
\textwidth 16cm
\textheight 21cm
\footskip 1.0cm


%~~~~~~~~~~~
% Section numbering
%~~~~~~~~~~~
\setcounter{secnumdepth}{0}  % don't number sections (stars not needed)

%~~~~~~~~~~~~
% Line numbering for equations
%~~~~~~~~~~~~
\usepackage{lineno}
\let\oldequation\equation
\let\oldendequation\endequation

\renewenvironment{equation}
{\linenomathNonumbers\oldequation}
{\oldendequation\endlinenomath}

%~~~~~~~~~~~
% Figures & Tables
%~~~~~~~~~~~
\usepackage{caption}
\usepackage{subcaption}

% Set the location of your figures
\graphicspath{{figs/}}  % in git location
%\graphicspath{{../../../../../../../Git/FR_n-prey-at-a-time/figs/}}  % in dropbox location


% Set the location of your tables
\newcommand*\inputtable[1]{\input{tables/#1.tex}}
%\newcommand*\inputtable[1]{\input{../../../../../../../Git/FR_n-prey-at-a-time/tables/#1.tex}}



% Uncomment following to place all figures at end
%\usepackage{caption}
%\captionsetup{labelsep=none,textformat=empty}
%\usepackage[figuresonly]{endfloat}
%\renewcommand{\listfigurename}{Figure captions}
%\usepackage{tocloft}
%\setlength{\cftfigindent}{0pt}
%\renewcommand{\cftfigpresnum}{Fig. }
%\addtolength{\cftfignumwidth}{1.5em}


%~~~~~~~~~~~
% Activative links
%~~~~~~~~~~~
\usepackage[hidelinks]{hyperref}

%~~~~~~~~~~~
% Citation setup
%~~~~~~~~~~~
% In-text citation styles
\usepackage[sort]{natbib}
\bibpunct[; ]{(}{)}{;}{a}{,}{;}

% Separate references and table of contents for the main text and appendix
\usepackage{appendix}
\usepackage{minitoc}

% Make the "Part I" text invisible
\renewcommand \thepart{}
\renewcommand \partname{}

% Automatic citiation count
\usepackage{totcount}
\newtotcounter{citnum}
\def\oldbibitem{} \let\oldbibitem=\bibitem
\def\bibitem{\stepcounter{citnum}\oldbibitem}

% Automatic figure and table count
\regtotcounter{figure}
\regtotcounter{table}

%~~~~~~~~~~~
% Table attributes
%~~~~~~~~~~~
% To force table caption above table
\usepackage{float}
\floatstyle{plaintop}
\restylefloat{table}

%  To rotate page to horizonal format for table
\usepackage{lscape}

% For multi-page tables
\usepackage{longtable}

%~~~~~~~~~~~
% Boxes
%~~~~~~~~~~~
% For making a multi-page box
\usepackage[most]{tcolorbox}


%~~~~~~~~~~~
% For non-numeric enumeration
%~~~~~~~~~~~
\usepackage{enumerate}

%-----------------------
% Line-referencing
%-----------------------
%  command to label changes
\newcommand{\LL}[1]{\linelabel{#1}}
%\newcommand{\LR}[1]{\lineref{#1}}

%~~~~~~~~~~~
% Environment for the abstract and cover page
%~~~~~~~~~~~
\newenvironment{myabstract}{\clearpage\section*{Abstract}}{\clearpage}
\newenvironment{mysignificance}{\clearpage\section*{Significance}}{\clearpage}
\newenvironment{cover}{\maketitle}{\clearpage}
\newenvironment{myfigure}[1]{\begin{figure}[#1]\centerline{Figure
			\arabic{figure}\medskip}}{\addtocounter{figure}{1}\end{figure}}


%~~~~~~~~~~~
% For annotation
%~~~~~~~~~~~
% For purposes of highlighting needed numbers
\newcommand{\redX}{\textcolor{red}{ X}\MN{XXX} }

% For commenting purposes (doesn't work in boxes!)
\long\def\authornote#1{%
	\leavevmode\unskip\raisebox{-3.5pt}{\rlap{$\scriptstyle\diamond$}}%
	\marginpar{\raggedright\hbadness=10000
		\def\baselinestretch{0.8}\tiny
		\it #1\par}}
\newcommand{\MN}[1]{\authornote{MN: #1}}
\newcommand{\KEC}[1]{\authornote{KEC: #1}}
% copy the last line and edit with co-author initials


%~~~~~~~~~~~
% Define convenience functions
%~~~~~~~~~~~
\DeclareMathOperator{\E}{\mathbb{E}} % For Expectation symbol
\renewcommand{\L}{\mathcal{L}} % likelihood symbol
\newcommand{\lL}{\ln \mathcal{L}} % log-likelihood symbol

%%%%%%%%%%%%%%%%%%%%%%%%%%%%%%%%%%%%%%%%%%%%
%%%%%%%%%%%%%%%%%%%%%%%%%%%%%%%%%%%%%%%%%%%%
%%%%%%%%%%%%%%%%%%%%%%%%%%%%%%%%%%%%%%%%%%%%

% Include your paper's title here
\title{Title}

\author
{Mark Novak$^1$\\\\
	\small{$^1$Department of Integrative Biology, Oregon State University,}\\
	\small{Corvallis, Oregon, 97331 USA}\\
	\\
}

% Include the date command, but leave its argument blank.
\date{}

%%%%%%%%%%%%%%%%%%%%%%%%%%%%%%%%
%%%%%%%%%% END OF PREAMBLE   %%%%%%%%%%
%%%%%%%%%%%%%%%%%%%%%%%%%%%%%%%%

\begin{document}
	
\maketitle

% ---------------------------------------
% COVER PAGE
% ---------------------------------------
\begin{cover}
	
	\centerline{{\sc Running Title:} XXXXX}
	
	\bigskip
	
	\begin{center}
		\begin{tabular}{ll}
			\hline 
			\\
			{\bf Article type}							 & Original Research             \\
			{\bf Words in abstract}					& XXX                                 \\
			{\bf Words in manuscript}       	 & XXX                               \\
			{\bf Number of references}			& \total{citnum}\                 \\
			{\bf Number of figures}					& \total{figure}                     \\
			{\bf Number of tables}					& \total{table}\                      \\
			{\bf Number of boxes}				   & XXX\                                  \\
			{\bf Corresponding author}		   & Mark Novak                    \\
			{\bf Phone} 								   & (541) 737-3610               \\
			{\bf Fax} 									  	  & (541)  737-3120              \\
			{\bf Email} 									& mark.novak@oregonstate.edu     \\
			\\
			\hline
		\end{tabular}
	\end{center}
	
\end{cover}


% --------------------------------------------------------------------------------
% --------------------------------------------------------------------------------
% SUMMARY OF REMAINING TO-DO'S
% --------------------------------------------------------------------------------
% --------------------------------------------------------------------------------
\baselineskip9pt
\noindent
\textbf{To do \& to discuss:}
\begin{itemize}
	\item{General}
	\begin{itemize}
		\item 
	\end{itemize}
	\item Results
	\begin{itemize}
		\item
	\end{itemize}
	\item Discussion
	\begin{itemize}
		\item 
	\end{itemize}
	\item SOM
	\begin{itemize}
		\item 
	\end{itemize}
\end{itemize}

\noindent
\textbf{Potential reviewers:}\\\\
	

\clearpage

\baselineskip12pt
%\tableofcontents



% ---------------------------------------
% ABSTRACT 
% ---------------------------------------

\baselineskip20pt
%\setpagewiselinenumbers
\linenumbers

\begin{myabstract}
	
	\lipsum[66]
		
	\noindent\textbf{Keywords:} 
		\emph{
			predator-prey
		}
\end{myabstract}



%\begin{mysignificance}
%	
%	\lipsum[66]
%	
%\end{mysignificance}


% -----------------------------------------------------------------------------
% -----------------------------------------------------------------------------
% BODY OF THE PAPER
% -----------------------------------------------------------------------------
% -----------------------------------------------------------------------------
\clearpage

\baselineskip30pt

% -----------------------------------------------------------------------------
% INTRODUCTION
% -----------------------------------------------------------------------------

\section{Introduction}

When using natbib, this is how to cite other papers \citep{kopka2004glt}.
Of course, using natbib, you can also cite \citet{lamport1989ldp} like this.
\MN{This is a comment}
\LL{P1} % hidden label to line number

%The line number was \LR{P1}.j


% -----------------------------------------------------------------------------
% Boxes
% -----------------------------------------------------------------------------

\begin{tcolorbox}[breakable, enhanced, before 
upper={\parindent15pt},boxsep=2mm]
	% set boxsep=0mm if not using line numbers
		\begin{internallinenumbers}
	\noindent
	\textbf{Box 1. Box title.}
	\lipsum[75]
	
	
		\end{internallinenumbers}
\end{tcolorbox}
%\linenumbers
\clearpage


% -----------------------------------------------------------------------------
% METHODS
% -----------------------------------------------------------------------------

\section{Methods}
\lipsum[66]

\subsection{Study system}
A subsection is initiated like this.
\lipsum[75]



% -----------------------------------------------------------------------------
% RESULTS
% -----------------------------------------------------------------------------

\section{Results}
References a figure by referencing its label (Fig.~\ref{fig:fig1}).
You can include a table either by making the table within this tex file, or call 
from a separate file (e.g., an R-generated table) using the \texttt{input} function.
As long as it has a label, you can reference your table the same way you reference figures and sections (Table~\ref{tab:table1}).

%\input{exampleTable}j


\begin{figure}[!h]
	\centerline{\includegraphics*[width=0.5\textwidth]{exampleFig.pdf}}
	\caption{
		\baselineskip30pt
		This is my first awesome figure.
	}
	\label{fig:fig1}
\end{figure}


% -----------------------------------------------------------------------------
% DISCUSSION
% -----------------------------------------------------------------------------

\section{Discussion}

\lipsum[75]





% -----------------------------------------------------------------------------
% ACKNOWLEDGMENTS
% -----------------------------------------------------------------------------
\section{Acknowledgments}
We thank everyone we know

\section{Author contributions}
Both authors contributed equally.

\section{Code and data availability}
It's all on GitHub! \emph{url address here}


% -----------------------------------------------------------------------------
% -----------------------------------------------------------------------------
% REFERENCES
% -----------------------------------------------------------------------------
% -----------------------------------------------------------------------------

% specify the bibiolography style file (without its file extension)
\bibliographystyle{biblio/ecology_letters2}
% specify your bib file (without its file extension)
\bibliography{biblio/bibliography}



\end{document}
